\documentclass[10pt,a4paper]{amsart}
\usepackage[margin=19mm]{geometry}
\usepackage{amsmath,amssymb,mathtools}
\usepackage[scr=rsfs,cal=euler]{mathalfa}
\usepackage{xfrac}
\usepackage[unicode,hidelinks,bookmarks=false]{hyperref}
\newcommand{\ie}{i.e.\ }
\newcommand{\cf}{cf.\ }
\newcommand{\resp}{respectively\ }
\newcommand{\Subsection}{\subsection}
\providecommand{\nobreakdash}{\nobreak\hbox{-}}
\setlength{\emergencystretch}{2em}
\multlinegap0pt
\usepackage{bm}
\usepackage{bbm}
\usepackage{dsfont}
\usepackage{xspace}
\usepackage{cancel}
\usepackage{mathtools}
\usepackage{amsthm}
\makeatletter
\@ifundefined{proposition}{\newtheorem{proposition}{Proposition}}{}
\@ifundefined{theorem}{\newtheorem{theorem}[proposition]{Theorem}}{}
\@ifundefined{corollary}{\newtheorem{corollary}[proposition]{Corollary}}{}
\makeatother
\DeclareMathOperator{\SO}{SO}
\DeclareMathOperator{\Spec}{Spec}
\DeclareMathOperator{\Sym}{Sym}
\DeclareMathOperator{\ord}{ord}
\newcommand{\HH}{\mathrm{H}}
\newcommand{\R}{\mathbb{R}}
\newcommand{\Z}{\mathbb{Z}}
\newcommand{\bC}{\mathbb{C}}
\newcommand{\cO}{\mathcal{O}}
\newcommand{\height}{\mathrm{Ht}}
\renewcommand{\O}{\mathcal{O}}
\renewcommand{\P}{\mathbb{P}}
\title[Kummers, spinors, and heights]{Kummers, spinors, and heights}
\author{Jef Laga, Jack A. Thorne}
\date{Source: arXiv:2507.06865v1 (CC BY 4.0)}
\begin{document}
\maketitle

\noindent\emph{Source and licence.} Kummers, spinors, and heights. Version arXiv:2507.06865v1 (CC BY 4.0), distributed under the Creative Commons Attribution 4.0 International licence, as recorded in the arXiv licence metadata for this exact version and shown on the abstract page https://arxiv.org/abs/2507.06865v1. Full rights notice: licenses/arxiv-2507.06865-CC-BY-4.0.txt.\par\smallskip

\noindent\textit{Editorial scope.} Counted unit: the Theorem whose source label is \texttt{thm\_canonical\_height\_difference} in the article (source0.tex lines 1850--1858, source label \texttt{thm\_canonical\_height\_difference}), delivered together with the article's own complete original proof of it (source0.tex lines 1859--1900, one \texttt{proof} environment of 42 lines). No mathematics is written, repaired or re-derived: statement and proof are the article's own text line for line. Required context retained uncounted: prop\_Riemann\_theta\_relation\_as\_incidence\_relation (lines 1445--1454); prop\_q\_tildes\_inject\_to\_Kummer (lines 1525--1531); cor\_properties\_of\_T\_Heisenberg\_matrices (lines 1603--1612); theorem\_weierstrassmodelregular\_muv\_zero (lines 1805--1811). Every definition, display and label of those lines is retained unchanged. Omitted from this package and not redistributed: 1--1444, 1455--1524, 1532--1602, 1613--1804, 1812--1849, 1901--2146. Nothing inside a retained excerpt is omitted or altered: every retained body is delivered exactly as the article's lines are written, so no omission receipt of any kind accompanies this package. Citations: the retained excerpts carry no citation command at all, so no bibliography entry of the article is transcribed and no reference is redistributed. Reference adaptation: none was needed and none was made; every reference inside the delivered unit resolves inside it, so no unresolved reference or citation is produced. Printed slips kept literal: no printed word, symbol, hypothesis or number of the counted statement or of its proof was altered. Layout and class: the article's own class is \texttt{article} with packages including amsmath, amssymb, amsthm, xy, fontenc, inputenc, enumitem, graphicx, geometry, hyperref, tikz, tikz-cd, xcolor, stmaryrd; the wrapper is the lane's pinned \texttt{amsart} editorial wrapper at 10pt on A4, which restates the theorem-like environments the retained text needs and adds the article's own macro definitions that the retained text needs (\texttt{\textbackslash HH}, \texttt{\textbackslash O}, \texttt{\textbackslash P}, \texttt{\textbackslash R}, \texttt{\textbackslash SO}, \texttt{\textbackslash Spec}, \texttt{\textbackslash Sym}, \texttt{\textbackslash Z}, \texttt{\textbackslash bC}, \texttt{\textbackslash cO}, \texttt{\textbackslash height}, \texttt{\textbackslash ord}), with extra wrapper packages (bm, bbm, dsfont, xspace, cancel, mathtools). The delivered numbering is the unit's own. Assets: no retained span contains a figure environment, a graphics-inclusion command, a PDF-page inclusion, a table or a TikZ picture; no image file is copied into this package, the package declares no asset dependency and no image file is redistributed with it. Licence: version arXiv:2507.06865v1 of this article is distributed under the Creative Commons Attribution 4.0 International licence, as recorded in the arXiv licence metadata for this exact version and shown on the abstract page https://arxiv.org/abs/2507.06865v1. A full rights notice is delivered at licenses/arxiv-2507.06865-CC-BY-4.0.txt. Characters: every retained excerpt is pure ASCII, so the delivered engine, XeLaTeX with its default Latin Modern text font, reports no missing character.
\begin{proposition}\label{prop_Riemann_theta_relation_as_incidence_relation}
    Let $T \in J[2]$ be a 2-torsion point of Mumford degree $g$, and let ${\widetilde{T}}$ be a pre-image in $\mathcal{G}(2 \Theta)$. Then:
    \begin{enumerate}
        \item The bilinear form $\beta( \widetilde{T} \cdot x, y)$ is symmetric; we write $q_{\widetilde{T}}(x) = \beta( \widetilde{T} \cdot x,x) \in \Sym^2 S^\vee =\HH^0(\P(S), \cO_{\P(S)}(2))$ for the associated quadratic form.
        \item Let $f_{\widetilde{T}}(x) = \beta(\widetilde{T} \cdot \infty,x) \in S^\vee =\HH^0(
        \P(S), \cO_{\P(S)}(1))$. Then we have the identity
        \[ f_{\widetilde{T}}(\infty) \cdot \xi([2]^\ast \Psi^\ast f_{\widetilde{T}}) = \Psi^\ast(q_{\widetilde{T}})^2 \]
        in $\HH^0(J, \cO_J(2 \Theta)^{\otimes 4})$.
    \end{enumerate}
\end{proposition}

\begin{proposition}\label{prop_q_tildes_inject_to_Kummer}
    \begin{enumerate}
        \item The action of $S \Gamma$ on $(\Sym^2 S^\vee) \otimes N$ factors through $\rho : S \Gamma \to \SO(V)$. Under the induced action of $J[2]$, there is an isomorphism $(\Sym^2 S^\vee) \otimes N \cong \oplus_S k(e_2(S, -))$, where $e_2 : J[2] \times J[2] \to \mu_2$ is the Weil pairing, and $S$ ranges over the set of points $S \in J[2]$ such that the theta-characteristic $(g-1) P_\infty + S$ is even. 
        \item Choose for each $T \in J[2]$ of Mumford degree $g$ a pre-image $\widetilde{T} \in \mathcal{G}(2 \Theta)$. Then the elements $\beta(\widetilde{T} x,x) \in\HH^0(\P(S), \cO_{\P(S)}(2))$ are linearly independent, and span a complement to the kernel $I_K$ of the map
        \[ \Psi^\ast :\HH^0(\P(S), \cO_{\P(S)}(2)) \to\HH^0(J, \O_J(2 \Theta)^{\otimes 2}). \]
    \end{enumerate}
\end{proposition}

\begin{corollary}\label{cor_properties_of_T_Heisenberg_matrices}
    \begin{enumerate}
        \item  The map $J[2] \to \mathcal{G}(2 \Theta), T \mapsto M_T$, is defined over $k$.
        \item Let $T \in J[2]$. Then:
        \begin{enumerate} \item We have $M_T^2 = r_T$ and $N(M_T) = q(T) r_T$, where $q : J[2] \to \mu_2$ is the quadratic form defined by $q(T) = (-1)^{|I| / 2}$ if $T = [D_I]$ and $|I|$ is even.
        \item The coefficients of $M_T$ lie in $\Z[\omega_1, \dots, \omega_{2g+1}]$.
        \item If $T$ has Mumford degree $g$, then $\beta(M_T \infty, \infty) = 1$; otherwise, $\beta(M_T \infty, \infty) = 0$.
        \end{enumerate}
    \end{enumerate}
\end{corollary}

\begin{theorem}\label{theorem_weierstrassmodelregular_muv_zero}
    Let $v$ be a finite place of $k$ of odd residue characteristic.
    Let $\mathcal{C}\rightarrow \Spec(\O_{k_v})$ be the Weierstrass model of $C$ over $\O_{k_v}$, given by the same equation $y^2 = x^{2g+1}+ c_1x^{2g} +\cdots c_{2g+1}$.
    Suppose that $\mathcal{C}$ is a regular scheme. 
    Then $\log \max_i(|x_i|_v/|x_1|_v)$ is a N\'eron function for $2\Theta$.
    Consequently, $\mu_v\equiv 0$. 
\end{theorem}

\begin{theorem}\label{thm_canonical_height_difference}
    \begin{enumerate} \item Let $v$ be a finite place of $k$, and let $P \in J(k_v)$. Then: 
    \begin{enumerate}
        \item If $\ord_v(\Delta) \leq 1$ and $v$ has odd residue characteristic, then $\mu_v \equiv 0$.
        \item Without any assumptions on $v$ or $\ord_v(\Delta)$, we have $\mu_v(P) \geq  \frac{1}{3} \log |2^{4 g} \Delta|_v$. 
    \end{enumerate}
    \item Let $v$ be an infinite place of $k$, corresponding to an embedding $\sigma : k \to \bC$, and let $P \in J(k_v)$. Then there is a constant $c = c(g) \in \R$ (independent of $f$ and $k$) such that $\mu_v(P) \geq \frac{1}{3} \log |\Delta|_v - \frac{1}{6} g(7g+5) \log \height(\sigma(f)) + c$.
    \end{enumerate} 
\end{theorem}

\begin{proof}
    (1) Part (a) follows from Theorem \ref{theorem_weierstrassmodelregular_muv_zero}, once we verify that the assumptions imply that the Weierstrass model $\mathcal{C}/\O_{k_v}$ is regular. 
    Since the double cover $\mathcal{C} \rightarrow \P^1_{\O_{k_v}}$ is \'etale outside the ramification locus $\mathcal{W}\sqcup \{P_{\infty}\}$ and since $\mathcal{C}\rightarrow \Spec(\O_{k,v})$ is smooth in a neighbourhood of the section $P_{\infty}$, it suffices to prove that every point of $\mathcal{W}$ is a regular point of $\mathcal{C}$.
    Since $\mathcal{W}$ is cut out by the regular element $y$, it suffices to prove that $\mathcal{W}$ is itself regular.
    But $\mathcal{W} = \Spec\O_{k_v}[x]/(f(x))$ and since $\ord_v(\Delta)\leq 1$, $\O_{k_v}[x]/(f(x))$ is a maximal order in the \'etale $k_v$-algebra  $k_v[x]/(f(x))$.
    It follows that $\O(\mathcal{W})$ is regular, hence $\mathcal{W}$ is regular, hence $\mathcal{C}$ is regular, as claimed. (This argument shows that the condition $\ord_v(\Delta) \leq 1$ could be weakened to the condition that $\cO_{k_v} / (f(x))$ is maximal order.)

    We now consider part (b). In this case, we will show that $\varepsilon_v(P) \geq \log | 2^{4g} \Delta |_v$ for all $P\in  J(k_v)$. We are free to enlarge $k$, so we can assume without loss of generality that all roots of $f$ lie in $k$. We consider the short exact sequence of $J[2]$-modules (cf. Proposition \ref{prop_q_tildes_inject_to_Kummer})
    \[ 0 \to I_K \otimes N \to (\Sym^2 S^\vee) \otimes N \to\HH^0(J, \O_J(2\Theta)^{\otimes 2}) \otimes N. \]
    According to Proposition \ref{prop_q_tildes_inject_to_Kummer}, there is an isomorphism of $J[2]$-modules 
    \[ (\Sym^2 S^\vee) \otimes N \cong \oplus_S k(e_2(S, -)), \]
    where the sum runs over the set of all $S \in J[2]$ such that $(g-1) P_\infty + S$ is an even theta-characteristic. Given such an $S$, let $\pi_S$ denote the projection to the $e_2(S, -)$-eigenspace; thus we have the formula
    \[ \pi_S = \frac{1}{2^{2g}} \sum_{T \in J[2]} e_2(S, T) N(M_T) M_T^{-1} = \frac{1}{2^{2g}} \sum_{T \in J[2]} e_2(S, T) N(M_T)^{-1} M_T = \frac{1}{2^{2g}} \sum_{T \in J[2]} e_2(S, T) (\pm r_{T})^{-1} M_T, \]
    where we have used Corollary \ref{cor_properties_of_T_Heisenberg_matrices}. Proposition \ref{prop_q_tildes_inject_to_Kummer} shows that there are constants $\lambda_{S, j} \in k$ such that if $S$ has Mumford degree $g$, we have
    \[ \pi_S(x_j^2) = \lambda_{S, j} q_{M_S}. \]
    Evaluating both sides at the point $\infty$, we find
    \[ \lambda_{S, j} =\frac{1}{2^{2g}} \sum_{T \in J[2]} e_2(S, T) (\pm r_{T}^{-1}) (M_T \infty)_j^2, \]
    hence $|\lambda_{S, j}|_v \leq | 2^{4g} \Delta |_v^{-1/2}$. 
    
    On the other hand, Proposition \ref{prop_Riemann_theta_relation_as_incidence_relation} shows that if $z \in \widetilde{K}(k_v)$, $w = \delta(z)$, then
    \[ \beta(M_S \infty, w) = q_{M_S}(z)^2. \]
    If $S$ is non-generic, then $\pi_S(x_j^2)$ lies in $I_K$, so vanishes at $z$. We find that
    \[ z_j^2 = \sum_{S \text{ generic}} \pi_S(x_j^2)|_{x = z} = \sum_{S \text{ generic}} \lambda_{S, j} q_{M_S}(z), \]
    hence 
    \[ \max_j |z_j|^2 \leq | 2^{4g} \Delta |_v^{-1/2} \max_S |q_{M_S}(z)|_v, \]
    hence
    \[ \max_j |z_j|^4 \leq |2^{4g} \Delta|_v^{-1} \max_S |\beta( M_S \infty,w)|_v \leq |2^{4g} \Delta|_v^{-1} \max_j |w_j|_v. \]
    Taking logarithms and re-arranging, this says exactly that, if $P \in J(k_v)$ is a point mapping to $[z] \in K(k_v)$, then
    \[ \varepsilon_v(P) \geq \log |2^{4g} \Delta|_v. \]
    This completes the proof of part (b).

    (2) When $v$ is an infinite place, we can argue in a similar way to Part 1(b), using the usual triangle inequality in lieu of the non-archimedean one.
    Identifying each term with its image under $\sigma$, and therefore dropping the subscript $v$ on each absolute value, we obtain
    \[ | \lambda_{S, j} | \ll \max_{T \in J[2]} \max_j | (M_T \infty)_j |^2 |\delta_T \delta_{T^c}| |\Delta|^{-1/2}.  \]
    Analysis of the behaviour under scaling shows that $| (M_T \infty)_j | \ll \height(\sigma(f))^{g(g+1)/2}$. Since $|\delta_T| \ll \height(\sigma(f))^{g(g-1)/2}$, we find 
    \[ | \lambda_{S, j} | \ll \height(\sigma(f))^{g(3g+2)/2} |\Delta|^{-1/2}. \]
    It follows that
    \[ \max_j | z_j |^4 \ll |\Delta|^{-1} \height(\sigma(f))^{g(3g+2) + g(g+1)/2} \max_j | w_j |. \]
    Re-arranging now gives
    \[ \varepsilon_v(P) \geq \log | \Delta | - \frac{1}{2} g(7g+5) \log \height(\sigma(f))  + O_g(1), \]
    as required. 
\end{proof}

\end{document}
